% =============================================================================
% Claude for Financial Services — 2-page A4 handout
% Audience leave-behind. Printable. Single column for readability.
%
% Compile:
%   latexmk -pdf -xelatex claude-financial-services-handout.tex
% =============================================================================
\documentclass[10pt,a4paper]{article}

% --- Geometry & spacing ------------------------------------------------------
\usepackage[a4paper,margin=15mm,top=12mm,bottom=12mm]{geometry}
\usepackage{parskip}
\setlength{\parskip}{4pt plus 1pt}
\pagestyle{empty}

% --- Encoding & fonts --------------------------------------------------------
\usepackage[english]{babel}
\usepackage{microtype}
\usepackage{fontspec}
\usepackage{amssymb}     % \checkmark

% --- Colour palette (matches the decks) --------------------------------------
\usepackage{xcolor}
\definecolor{accentpurple}{HTML}{7C3AED}
\definecolor{darkgray}{HTML}{1F2937}
\definecolor{lightgray}{HTML}{F3F4F6}
\definecolor{successgreen}{HTML}{059669}
\definecolor{warningred}{HTML}{DC2626}

% --- Boxes for callouts ------------------------------------------------------
\usepackage[most]{tcolorbox}
\tcbset{
  boxrule=0pt,
  arc=2pt,
  left=6pt,right=6pt,top=4pt,bottom=4pt,
  boxsep=2pt,
  fonttitle=\bfseries\small,
}
\newtcolorbox{callout}[1][]{colback=lightgray, colframe=lightgray, #1}
\newtcolorbox{accentbox}[1][]{colback=lightgray, colframe=accentpurple, boxrule=1.2pt, leftrule=3pt, #1}

% --- Tables ------------------------------------------------------------------
\usepackage{booktabs}
\usepackage{array}
\usepackage{tabularx}

% --- Section headers ---------------------------------------------------------
\usepackage{titlesec}
\titleformat{\section}
  {\color{accentpurple}\bfseries\large}
  {}{0pt}{}
  [\vspace{-6pt}\noindent\color{accentpurple}\rule{\textwidth}{0.6pt}]
\titlespacing*{\section}{0pt}{8pt}{4pt}
\titleformat{\subsection}
  {\color{darkgray}\bfseries\normalsize}
  {}{0pt}{}
\titlespacing*{\subsection}{0pt}{4pt}{2pt}

% --- Helpers -----------------------------------------------------------------
\newcommand{\good}{\textcolor{successgreen}{\textbf{\checkmark}}}
\newcommand{\bad}{\textcolor{warningred}{\textbf{\texttimes}}}
\newcommand{\hl}[1]{\textbf{\textcolor{accentpurple}{#1}}}

% --- Tighter list spacing ----------------------------------------------------
\usepackage{enumitem}
\setlist{nosep, leftmargin=*, topsep=2pt, itemsep=1pt}

% --- Hyperlinks --------------------------------------------------------------
\usepackage{hyperref}
\hypersetup{hidelinks, pdftitle={Claude for Financial Services — handout}}

% =============================================================================
\begin{document}

% --- Title bar ---------------------------------------------------------------
\noindent
\begin{tcolorbox}[colback=darkgray, colframe=darkgray, fonttitle=\Large\bfseries\color{white}]
  \color{white}\textbf{\Large Claude for Financial Services} \hfill \small Internal briefing pack \quad\textbar\quad June 2026 \\[2pt]
  {\small A printable summary of the May 2026 launch, where we stand with it, and what to say when people ask. Drill at your desk.}
\end{tcolorbox}

% =============================================================================
\section*{The headline}

On \hl{5 May 2026}, Anthropic launched Claude for Financial Services:

\begin{itemize}
  \item \textbf{10 ready-to-run agent templates} for high-volume FS workflows
  \item \textbf{Microsoft 365 integration} — Excel, PowerPoint, Word, Outlook
  \item \textbf{MCP-based connectors} to $\sim$11 financial data providers (Daloopa, S\&P, Moody's, FactSet, LSEG, PitchBook, $\ldots$)
  \item \textbf{Enterprise controls} — SSO (SAML/OIDC), SCIM, audit logs, ISO/IEC 42001:2023
\end{itemize}

\noindent Underlying model: \hl{Claude Opus 4.7}. Benchmark: \hl{64.37\%} on Vals AI Finance Agent (finance-specific eval, current best-in-class).

% =============================================================================
\section*{The 10 agents}

\begin{tabularx}{\textwidth}{@{}p{0.42\textwidth} X@{}}
\toprule
\textbf{Research \& Client Coverage} & \textbf{Finance \& Operations} \\
\midrule
\textbf{Pitch Builder} — target lists, comps, pitchbooks & \textbf{Valuation Reviewer} — vs comparables \& firm standards \\
\textbf{Meeting Preparer} — client/counterparty briefs & \textbf{GL Reconciler} — reconciliation \& NAV \\
\textbf{Earnings Reviewer} — transcripts → model updates & \textbf{Month-End Closer} — close checklist + reports \\
\textbf{Model Builder} — models from filings + feeds & \textbf{Statement Auditor} — consistency \& audit-readiness \\
\textbf{Market Researcher} — sector/issuer intel & \textbf{KYC Screener} — entity files + escalations \\
\bottomrule
\end{tabularx}

\smallskip
\noindent\textit{\small Each template is a bundle of three layers: markdown \textbf{skills} (how a senior banker does this), MCP \textbf{connectors} (scoped data access), and addressable \textbf{subagents}. Inspectable, not a black box.}

% =============================================================================
\section*{How it ships --- three access paths}

\begin{tabularx}{\textwidth}{@{}lXl@{}}
\toprule
\textbf{Path} & \textbf{What you get} & \textbf{Requires} \\
\midrule
A. M365 + Cowork  & Claude inside Excel/PPT/Word/Outlook + the 10 FS agents     & Claude Team / Enterprise + M365 admin install \\
B. Claude Code    & Agents from terminal / scripted context                     & Claude Code subscription \\
C. Managed Agents & Programmatic, autonomous, behind your workflow engine       & Anthropic API access (org) + engineering \\
\bottomrule
\end{tabularx}

\smallskip
\noindent\textit{\small Path A note: the M365 add-in can be deployed against your firm's own cloud (Vertex AI, Bedrock, internal LLM gateway) via the \texttt{claude-for-msft-365-install} Claude Code plugin — traffic then stays in your cloud, compliance posture is whatever you already have for that cloud.}

% =============================================================================
\section*{Where we stand}

\begin{tabularx}{\textwidth}{@{}X X@{}}
\toprule
\textbf{Have} & \textbf{Don't have} \\
\midrule
\good\ Microsoft 365 Copilot           & \bad\ Claude.ai (consumer) \\
\good\ GitHub Copilot                  & \bad\ Claude Team / Enterprise \\
                                       & \bad\ Claude for M365 (Excel/PPT/Word/Outlook plugin) \\
                                       & \bad\ Claude Cowork \\
                                       & \bad\ Claude Code \\
                                       & \bad\ Anthropic API (org-level) \\
\bottomrule
\end{tabularx}

\smallskip
\noindent\hl{No path} from M365 Copilot or GitHub Copilot to the Claude FS agents. The two product families are entirely separate.

% --- footer page 1 -----------------------------------------------------------
\vfill
\noindent{\footnotesize\color{darkgray}\textit{Verified 2026-06-08 against \texttt{anthropic.com/news/finance-agents}, \texttt{claude.com/solutions/financial-services}, and \texttt{github.com/anthropics/financial-services}. \hfill Page 1 of 2}}

\newpage

% =============================================================================
\section*{What to tell colleagues}

\begin{accentbox}
\textit{``Yes --- Anthropic launched Claude for Financial Services in May. It's a strong product for our type of work, but we don't have a sanctioned channel to use it: our office AI tooling today is Microsoft 365 Copilot and GitHub Copilot, and Claude requires a separate subscription we haven't acquired.''}
\end{accentbox}

\smallskip
\noindent If they push --- \textit{``can I just use my personal account?''}:

\begin{accentbox}
\textit{``Not on office data, no --- that would be a data-handling issue regardless of how useful the output is. Personal Claude is fine for personal exploration; office work goes through office-sanctioned tools.''}
\end{accentbox}

% =============================================================================
\section*{Personal accounts --- the boundary}

\begin{callout}
A personal Claude subscription is \hl{not} a sanctioned channel for office work. No DPA with our firm; no audit logging into our compliance posture; default Anthropic retention (not enterprise zero-retention). \textbf{Personal Claude is for personal exploration. Office work goes through office-sanctioned tools.} True regardless of how useful the output looks.
\end{callout}

% =============================================================================
\section*{FAQ}

\footnotesize
\begin{description}[leftmargin=0pt, labelindent=0pt, style=unboxed, font=\bfseries\color{darkgray}]
\item[Q. So can we get it?]
A: That's a procurement \& compliance conversation, not this one. Worth raising with whoever owns AI tooling decisions.

\item[Q. Isn't M365 Copilot enough?]
A: For a large slice, yes. Where it isn't: multi-doc synthesis, citation discipline, FS-specific connectors. Whether that gap matters is the procurement conversation.

\item[Q. Isn't this just the same model with a coat of paint?]
A: Partly fair. Underlying model is Claude Opus 4.7 (same as raw API). What's added: bundled FS templates, FS data connectors, M365 plugins, enterprise controls. Four things not trivial to assemble yourself.

\item[Q. Personal accounts on personal time?]
A: Personal data, personal projects --- fine. Office data through office tooling. Same rule as any external AI service.

\item[Q. Should we switch to Claude for Excel?]
A: On the FMW benchmark, Claude leads at 83\%. But that's a snapshot; Copilot is improving. ``Should we switch'' is the procurement conversation. Today we use what we have.

\item[Q. Can we run the demo template on our data?]
A: No. Public data only, in a personal Cowork instance, for demonstration. European banks, Microsoft, Apple --- fine. Office data --- not appropriate.

\item[Q. Couldn't we just clone the GitHub repo and self-host?]
A: The repo is Apache 2.0, but the orchestrator is Claude itself, which only runs via paid Anthropic surfaces (or BYOC on Vertex/Bedrock). Plus: data connectors require paid data subscriptions; templates are generic (firm-specific tuning is the 80\%). Free to read; not free to run.

\item[Q. Big Four --- Deloitte, PwC, KPMG, EY?]
A: All four announced internal Claude rollouts: PwC and KPMG in mid-May, Deloitte in late 2025, EY in late-May. Whether they're packaging FS consulting around it externally is less clear --- partnership conversation, not a tooling one.

\item[Q. Isn't sending Excel data to an external AI service the dealbreaker?]
A: Out of the box, that's where traffic goes. But the M365 add-in can be deployed via the \texttt{claude-for-msft-365-install} plugin against your firm's own cloud (Vertex / Bedrock / internal gateway). Data path becomes Excel → your cloud → Claude on your cloud → back.

\item[Q. I think we should buy it.]
A: Good. Right venue is whoever owns AI tooling decisions here. This pack is the briefing they'll need.
\end{description}
\normalsize

% =============================================================================
\section*{Further reading}

\subsection*{The team folder}
\footnotesize
\begin{itemize}
  \item \texttt{01-what-it-is.md} --- the offering in detail
  \item \texttt{02-how-it-ships.md} --- subscription tiers, deployment paths, BYOC routing
  \item \texttt{03-our-access-situation.md} --- the honest articulation of where we stand
  \item \texttt{04-what-we-can-do-today.md} --- pragmatic guidance with M365 + GitHub Copilot
  \item \texttt{demos/01-sector-overview.md} --- prompt template to run the Market Researcher agent on a public sector (e.g.\ European banks) in personal Cowork
  \item \texttt{appendix-sources.md} --- every claim, sourced; videos; what we did NOT verify
\end{itemize}

\subsection*{Videos (May 2026 launch)}
\begin{itemize}
  \item \textbf{The Briefing: Financial Services} (Anthropic event, 5 May 2026, 1.5h) --- the launch event itself.\\
        \texttt{anthropic.com/events/the-briefing-financial-services-virtual-event}
  \item \textbf{Claude for FS: Putting agents to work} (webinar, 7 May 2026, with demos).\\
        \texttt{anthropic.com/webinars/claude-for-financial-services-putting-agents-to-work}
  \item \textbf{Claude for FS teams --- on-demand recording} (59 min, Nick Lynn + Owen Scully) --- best single resource.\\
        \texttt{anthropic.ondemand.goldcast.io/on-demand/74ae0fc7-f591-442f-8fe1-244f409c1fb5}
\end{itemize}

\subsection*{Open-source repo (Apache 2.0)}
\begin{itemize}
  \item \texttt{github.com/anthropics/financial-services} --- skills, agents, MCP connectors, the \texttt{claude-for-msft-365-install} BYOC plugin. Free to read; running anything still requires paid Anthropic surfaces or BYOC on Vertex / Bedrock.
\end{itemize}
\normalsize

% --- footer page 2 -----------------------------------------------------------
\vfill
\noindent{\footnotesize\color{darkgray}\textit{Verified 2026-06-08. \hfill Page 2 of 2}}

\end{document}
